% ============================================================
% f2015ahw_q17_histogram.tex
% 2015 安徽卷（文科）第 17 题 · 频率分布直方图 —— 清单 F04 / §4
%
% 来源：.work/phase3-b12/figure-inventory.md §4（源 PDF 内容流逐 path 实测真页 pt）
%       源件 .work/phase3-b12/source/2015安徽文.pdf p1
%
% 坐标换算：画布 [x=1pt, y=1pt]（单位 TeX pt）。局部坐标 = 真页 pt − 原点
%   O = (866.533, 257.207)（清单 §4 几何 bbox 左下角）。真页 pt → 局部 pt 为
%   1:1 平移、不缩放——清单 §0.3 的 0.999622 是 pdftocairo 渲染件的 SVG 观测口径换算，
%   只用于与渲染件对读，不入绘图坐标。渲染层面 1 TeX pt = 0.996264 PDF pt（72/72.27），
%   整图统一 0.37%，属排版口径而非坐标换算误差。
%
% 线宽（清单 §0.4 本卷实测分档）：坐标轴、柱体、虚线导引线 0.398pt；纵轴分式横线 0.42pt
% 虚线：[2.989 2.989] 相位 0（5 条导引线）
% 箭头：arrows.meta Triangle 长 3.985pt、宽 2.988pt（清单 §6.3 全卷 31 枚同规格）；
%   箭尖落在清单尖端坐标；线端由 TikZ 按箭长回缩（源内线端再多深入 0.398pt = 一个
%   线宽，被填充三角完全覆盖，不可见）。
% 标签：字号 absolute 10.5pt（汉字 STSong 10.5、刻度数字 CMR/LMRoman 10.5、纵轴 a 为 CMMI 10.5 斜体）；一律 anchor=base 落在清单实测基线上——
%   清单给的是字框中心：汉字字框 = 全角 em 方框（基线 = 中心 − 0.36em），拉丁/数学字框
%   = 字体 ascent/descent 方框（基线 = 中心 − 0.25em）；两式均与清单另给的基线自洽。
%
% 照源复刻的存疑项：
%    - 6 柱（不是 8 柱；8 是横轴刻度个数），相邻柱各自成矩形、共用边两处相差 0.001pt（B-S3b），
%      照源逐柱独立画，不合并成连续条带、不改柱宽。
%    - a 柱顶 295.434 与其等高虚线导引线 295.448 相差 0.014pt（B-S3）——源内不一致，照源不对齐。
%    - 0→40 的轴长被压缩为 19.843pt（1 格，正常 4 格），其间以 3 峰 3 谷 zigzag 占位（B-S10）；
%      柱 [40,50) 的左边界 913.994 与刻度「40」文字中心重合，照源。
%    - 断轴 zigzag 与 x 轴主干在源中为同一条 14 节点折线（#119），本图按折线照录。
% ============================================================
\begin{tikzpicture}[x=1pt, y=1pt,
  every node/.style={inner sep=0, font=\fontsize{10.5pt}{12.6pt}\selectfont},
  line cap=butt, line join=miter,
  arr/.style={-{Triangle[length=3.985pt, width=2.988pt]}},
  dashline/.style={line width=0.398pt, dash pattern=on 2.989pt off 2.989pt}]

  % ---- 坐标轴：单条 14 节点折线（含断轴 zigzag），y 轴 145.251pt、x 轴主干 145.233pt ----
  \draw[line width=0.398pt] (27.618,157.980) -- (27.618,12.729) -- (31.586,12.729) -- (32.579,14.004) -- (34.563,11.453) -- (35.555,12.729) -- (36.547,14.004) -- (38.532,11.453) -- (39.524,12.728) -- (40.516,14.004) -- (42.501,11.453) -- (43.493,12.728) -- (47.461,12.728) -- (192.694,12.729);
  \draw[line width=0.398pt, arr] (27.618,157.980) -- (27.618,161.567);
  \draw[line width=0.398pt, arr] (192.694,12.729) -- (196.281,12.729);
  % 纵轴「频率/组距」分式横线（线宽 0.42pt）
  \draw[line width=0.42pt] (32.512,160.887) -- (53.512,160.887);

  % ---- 6 柱（各宽 19.843pt；逐柱独立矩形，共用边照源微差）----
  \draw (47.461,12.729) rectangle (67.304,29.728);
  \draw (67.303,12.729) rectangle (87.146,38.227);
  \draw (87.146,12.729) rectangle (106.989,106.288);
  \draw (106.989,12.729) rectangle (126.832,131.786);
  \draw (126.831,12.729) rectangle (146.674,106.288);
  \draw (146.674,12.729) rectangle (166.517,89.289);

  % ---- 5 条虚线导引线（起于纵轴 x=894.151，终于各自柱左上角）----
  \draw[dashline] (27.618,29.728) -- (47.461,29.728);
  \draw[dashline] (27.618,38.241) -- (67.303,38.241);
  \draw[dashline] (27.618,106.288) -- (126.831,106.288);
  \draw[dashline] (27.618,131.786) -- (106.989,131.786);
  \draw[dashline] (27.618,89.289) -- (146.674,89.289);

  % ---- 标签（清单 §4.4/§4.5；anchor=base 于实测基线）----
  \node[anchor=base] at (27.617,2.258) {$0$};
  \node[anchor=base] at (47.457,2.258) {$40$};
  \node[anchor=base] at (67.307,2.258) {$50$};
  \node[anchor=base] at (87.147,2.258) {$60$};
  \node[anchor=base] at (106.987,2.258) {$70$};
  \node[anchor=base] at (126.827,2.258) {$80$};
  \node[anchor=base] at (146.677,2.258) {$90$};
  \node[anchor=base] at (166.517,2.258) {$100$};
  \node[anchor=base] at (43.007,165.363) {频率};
  \node[anchor=base] at (43.007,151.063) {组距};
  \node[anchor=base] at (11.957,128.618) {$0.028$};
  \node[anchor=base] at (11.957,103.128) {$0.022$};
  \node[anchor=base] at (11.957,86.128) {$0.018$};
  \node[anchor=base] at (21.147,35.978) {$a$};
  \node[anchor=base] at (11.957,26.508) {$0.004$};
  \node[anchor=base] at (196.277,18.243) {分数};
\end{tikzpicture}
